% ==================================================
% CHAPTER 5: SECURITY AND PRIVACY
% ==================================================

\chapter{Security and Privacy}
\label{ch:security}

MemoryLane stores what people read, which can be personal. Its main security
property is \textbf{isolation}: nobody can read, search, summarise or delete
another person's pages. Its main privacy property is \textbf{restraint}: it
does not collect what it should not.

\section{Authentication}

Sign-in is handled by Amazon Cognito \cite{cognito}, with Google or email
and password; MemoryLane never handles a password. The web app and the
extension use the OAuth 2.0 authorization-code flow \cite{rfc6749} with PKCE
\cite{rfc7636} and a public client, so no client secret ships to the
browser.

\begin{figure}[!htbp]
\centering
\begin{tikzpicture}[x=3.05cm, y=0.5cm]
  \foreach \i/\t in {0/{User}, 1/{:Web app}, 2/{:Amazon\\Cognito}, 3/{:FastAPI}, 4/{:accounts.db}} {
    \node[lifehead] (h\i) at (\i,0) {\t};
    \draw[life] (h\i.south) -- (\i,-19.0);
  }
  \newcommand{\m}[4]{\draw[msg] (#1,#3) -- node[mlabel, fill=white] {#4} (#2,#3);}
  \newcommand{\rt}[4]{\draw[ret] (#1,#3) -- node[mlabel, fill=white] {#4} (#2,#3);}
  \newcommand{\self}[3]{\draw[msg] (#1,#2) -- ++(0.1,0) -- ++(0,-0.75) node[font=\scriptsize, right, pos=0.5, align=left, fill=white, inner sep=1pt] {#3} -- ++(-0.1,0);}

  \m{0}{1}{-1.6}{Sign in}
  \self{1}{-2.2}{make verifier and\\SHA-256 challenge}
  \m{1}{2}{-3.9}{/oauth2/authorize (challenge)}
  \m{2}{0}{-4.8}{hosted sign-in (Google or email)}
  \m{0}{2}{-5.7}{credentials}
  \rt{2}{1}{-6.6}{redirect with one-time code}
  \m{1}{2}{-7.5}{/oauth2/token (code + verifier)}
  \rt{2}{1}{-8.4}{access + refresh tokens}
  \m{1}{3}{-9.5}{API request, Bearer access token}
  \self{3}{-10.1}{verify RS256, exp,\\iss, token\_use, client\_id}
  \m{3}{4}{-11.8}{library for this sub?}
  % alt frame
  \draw (2.55,-12.4) rectangle (4.35,-17.6);
  \node[draw, fill=white, font=\tiny, anchor=north west, inner sep=1.5pt] at (2.55,-12.4) {alt};
  \node[font=\tiny, anchor=west] at (2.6,-12.95) {[known sign-in]};
  \rt{4}{3}{-13.5}{library id}
  \draw[densely dashed] (2.55,-14.1) -- (4.35,-14.1);
  \node[font=\tiny, anchor=east] at (4.3,-14.55) {[first time]};
  \m{3}{2}{-15.2}{/oauth2/userInfo (user's token)}
  \rt{2}{3}{-16.0}{email, email\_verified}
  \m{3}{4}{-16.9}{link sub to library of verified email}
  \rt{3}{1}{-18.4}{response (only that library's data)}
\end{tikzpicture}

\caption{Sequence diagram: sign-in (authorization code with PKCE) and the first API request}
\label{fig:seq-signin}
\end{figure}

The access token is kept only in memory and the refresh token in local
storage, so a reload stays signed in. Because a script on the page could read
that token, the app never renders user-supplied HTML, and Caddy sends a
strict Content Security Policy. The session ends only when Cognito actually
refuses the refresh token; a refresh cancelled by a reload, a network error
or throttling keeps the user signed in.

\section{One Library per Person}

Cognito gives a Google sign-in and an email sign-in different identifiers
even for the same address. The backend links them by \emph{verified} email
(ADR~0005): the first time an identifier is seen, it asks Cognito for the
email using the user's own token, so the browser cannot fake it. An
unverified sign-in keeps a separate library until it is verified, and its
pages are then moved across in the background. Links are re-checked daily,
and if Cognito cannot be reached for a new identifier the request fails
rather than guessing.

\section{Data Isolation}

Isolation is enforced in the data-access layer, not left to each endpoint:

\begin{itemize}
  \item The Solr client adds the caller's library filter to \emph{every}
  query, scan, KNN search and delete, and refuses to run without a valid id.
  \item KNN uses the filter as a pre-filter, so other users' vectors are
  never even candidates.
  \item The user id is part of every page id, so copies never collide.
  \item A request for another user's page returns \texttt{404}, exactly like
  a missing page, so ids cannot be probed.
  \item No endpoint accepts a user id from the client.
\end{itemize}

\section{Threat Model}

Table~\ref{tab:threats} relates the main threats, several from the OWASP
Top~10 \cite{owasp2021}, to their mitigations.

\begin{table}[!htbp]
\centering
\small
\begin{tabularx}{\textwidth}{p{0.27\textwidth} L}
\toprule
\textbf{Threat} & \textbf{Mitigation} \\
\midrule
Reading another user's pages (broken access control) & Library filter in the Solr client; pre-filtered KNN; 404 for foreign ids; isolation tests \\
Forged or replayed tokens & RS256 verification against Cognito's keys; expiry, issuer, token type and client checks \\
Stolen refresh token via injected script (XSS) & No user HTML rendered; strict CSP; tokens out of reach of pages in the extension \\
Query injection into Solr & Field syntax disabled; values escaped; ids validated \\
Server-side request forgery & The backend never fetches a client-supplied URL \\
Abuse and resource exhaustion & Per-user rate limits; 5{,}000-page cap; request size limits \\
Sensitive pages captured & Address, reading and content rules; server re-check; one-click removal \\
Personal data in logs & Query strings dropped from access logs; Solr and Cognito request URLs not logged \\
Exposed internals & Only ports 80/443 open; Solr bound to localhost; backend unpublished \\
\bottomrule
\end{tabularx}
\caption{Threats and mitigations}
\label{tab:threats}
\end{table}

\section{Privacy}

Automatic capture is off by default and never captures private pages
(Figure~\ref{fig:swimlane}). Logs hold page ids, never the URLs people save
or what they search for, and a test checks it. Page text leaves the server
only when the user asks for a summary of that page, and users can delete any
page, or all their data, at any time.
